\documentclass[15pt]{sprawozdanie}
\usepackage{soul}

\class{Praca dyplomowa inżynierska}
\title{Opracowanie wirtualnego środowiska do symulacji dynamiki lotu bezzałogowych statków powietrznych}
\author{\textbf{inż. Wojciech Gajda} 304494\\\vspace{20pt}\textbf{Igor Faliszewski} 313223}
\instructor{dr inż. Paweł Kotowski}
\deadline{\today}

\graphicspath{{images/}}

\usepackage{amssymb}
\usepackage{amsmath}
\usepackage{polski}
\usepackage[utf8]{inputenc}
\usepackage{hyperref}
\usepackage{blindtext}
\usepackage{multicol}
\usepackage{wrapfig}
\usepackage{float}
\usepackage{enumitem}
\usepackage{xfrac}
\usepackage{caption}
\usepackage{subcaption}
\usepackage{booktabs}
\usepackage{wasysym}
\usepackage{xcolor}
\usepackage{listings}
\usepackage{pdfpages}
\usepackage{fontspec}
\usepackage{comment}
\setmainfont{Verdana}

\definecolor{mGreen}{rgb}{0,0.6,0}
\definecolor{mGray}{rgb}{0.5,0.5,0.5}
\definecolor{mPurple}{rgb}{0.58,0,0.82}
\definecolor{backgroundColour}{rgb}{0.95,0.95,0.92}

\lstdefinestyle{CStyle}{
    backgroundcolor=\color{backgroundColour},   
    commentstyle=\color{mGreen},
    keywordstyle=\color{magenta},
    numberstyle=\tiny\color{mGray},
    stringstyle=\color{mPurple},
    basicstyle=\footnotesize,
    breakatwhitespace=false,         
    breaklines=true,                 
    captionpos=b,                    
    keepspaces=true,                 
    numbers=left,                    
    numbersep=5pt,                  
    showspaces=false,                
    showstringspaces=false,
    showtabs=false,                  
    tabsize=2,
    language=C
}

\usepackage[
%backend=biber
%,style=alphabetic
%,sorting=none
]{biblatex}
\addbibresource{bibliografia.bib}
%\DeclareNameAlias{author}{last-first}

\DeclareCiteCommand{\supercite}[\mkbibsuperscript]
  {\iffieldundef{prenote}
     {}
     {\BibliographyWarning{Ignoring prenote argument}}%
   \iffieldundef{postnote}
     {}
     {\BibliographyWarning{Ignoring postnote argument}}}
  {\usebibmacro{citeindex}%
   \bibopenbracket\usebibmacro{cite}\bibclosebracket}
  {\supercitedelim}
  {}
  
  \DeclareLabelalphaTemplate{
  \labelelement{
    \field[final]{shorthand}
    \field{label}
    \field[strwidth=3,strside=left,ifnames=1]{labelname}
    \field[strwidth=1,strside=left,final]{labelname}
    \field{labeltitle}
  }
  \labelelement{
    \field[strwidth=2,strside=right]{year}
  }
}
  
\let\cite=\supercite

\renewcommand*{\figurename}{Rys.}

\usepackage{titlesec}
\titlelabel{\thetitle.\quad}

\usepackage{tikz}


\usetikzlibrary{matrix, ,backgrounds}
\usepackage{array}

\makeatletter
\tikzset{SWOT/.style={matrix of nodes,inner sep=0pt,row sep=0pt,column sep=0pt,
cells={nodes={anchor=center,inner sep=2pt}},
column 1/.style={nodes={rotate=90,minimum height=8mm}},
ampersand replacement=\&,
execute at end matrix={\begin{scope}[on background layer]
 \fill[black!10] (\tikz@fig@name.west|-\tikz@fig@name-2-2.north) rectangle 
  (\tikz@fig@name-\the\pgfmatrixcurrentrow-2.south west);
\end{scope}
\draw (\tikz@fig@name.west|-\tikz@fig@name-2-2.north) rectangle 
(\tikz@fig@name-\the\pgfmatrixcurrentrow-\the\pgfmatrixcurrentcolumn.south east)
 (\tikz@fig@name-1-2.north west) rectangle 
(\tikz@fig@name-\the\pgfmatrixcurrentrow-\the\pgfmatrixcurrentcolumn.south east)
(\tikz@fig@name-2-2.center|-\tikz@fig@name.north) --
 (\tikz@fig@name-2-2.center|-\tikz@fig@name.south)
foreach \XX in {2,...,\the\numexpr\pgfmatrixcurrentrow-1}
{(\tikz@fig@name-\XX-2.south-|\tikz@fig@name.west) --
(\tikz@fig@name-\XX-2.south-|\tikz@fig@name.east) };
}}}
\makeatother

\usepackage{tocloft}
\renewcommand\cftfigfont{\small}

\begin{document}
%%\maketitle
\includepdf[pages=-]{first_page.pdf}


\section*{Streszczenie}
W pracy opisano realizację systemu przeznaczonego do symulacji dynamiki lotu bezzałogowych statków powietrznych. System pozwala na prowadzenie symulacji lotu w czasie rzeczywistem, który dodatkowo jest prezentowany jest w postaci trójwymiarowej wizualizacji. W trakcie wykonywania lotu logowane są dane i mogą zostać wykorzystane do analizy lotu. Opracowany został uniwersalny model dynamiki pozwalający na swobodną konfigurację parametrów statku. Obejmuje on modyfikację właściwości mechanicznych, aerodynamicznych oraz konfigurację zespołów napędowych i~wpływu czynników zewnętrznych. Symulacja dynamiki rozszerzona została o system sterowania. System został zaprojektowany w sposóbu ułatwiający zmianę parameterów statków i symulacji, tworzenie nowych konfiguracji statków oraz tworzenie i strojenie systemów sterowania. W pracy zaprezentowane zostało działanie systemu dla przykładowych modeli latających: stałopłatowca, wielowirnikowca i rakiety.

\section*{Słowa kluczowe}

symulacja, grafika komputerowa 3D, bezzałogowy statek powietrzny, model dynamiki ruchu

\newpage

\section*{Abstract}

\section*{Keywords}

\newpage
\tableofcontents

\newpage



\color{red}

\section{Wstęp}

\subsection{O projekcie}

Symulacje komputerowe dynamiki ruchu stanowią użyteczne narzędzie w pracach inżynierskich. Pozwalają na analizę poprawności działania układu mechanicznego przed jego wyprodukowaniem. W szczególności w zagadnieniu jakim jest projektowanie bezzałogowych statków powietrznych, zastosowanie symulacji pozwala zminimalizować koszty wytworzenia poprawnie działającego systemu.

\cite{rust}


\subsection{Cel projektu}

Celem niniejszej pracy jest opracowanie wirtualnego środowiska do symulacji dynamiki lotu bezzałogowych statków powietrznych. System implementuje podstawowy model dynamiki statków powietrznych wyposażonych w silniki rotorowe, silniki odrzutowe, powierzchnie nośne i powierzchnie sterowe. Pozwala na przeprowadzenie lotu symulowanym obiektem ktorego parametery określane są przez konfiguracje. System dzieli sie na serwer i aplikacje kliencką. Róznorodność modułów pozwala na realizacje róznych scenariusz (strzał, upuszczenie ładunku, kolizje, wpływ warunków srodowiskowych etc.)

\subsection{Przegląd istniejących rozwiązań}

Historia symulatorów lotu sięga lat 30 XX wieku. Pierwotnie zastosowanie symulatorów sprowadzało się do szkolenia pilotów cywilnych i wojskowych. W znanej obecnie formie kompletne symulatory lotu stanowią rozbudowane systemy integrujące wysokiej klasy oprogramowanie z peryferiami mającymi wierne odwzorowanie kokpitu kierowanej maszyny. Symulatory wykorzystywane do treningu pilotów podlegają rygorystycznym regulacją prawnym i na ogół ich zadaniem jest odwzorowanie jednej konkretnej maszyny. Równolegle uproszczone wersje symulatorów zaczeły zyskiwać popularność w zastosowaniu cywilnym, jako element rozrywki. W szczególności gry komputerowe związane z lotem bardzo często poświęcały zgodność z model rzeczywistym na rzecz lepszych odczuć użytkownika.

Na rynku dostępnych jest wiele środowisk symulacyjnych o różnym stopniu szczegółowości. Pełne systemy lotu stanowią produkt komercyjny projektowany na indywidualne zamówienie. Do najpopularniejszych dostępnych systemów sprzedawanych jako zamkniete oprogramowanie należą m. in.:

\begin{itemize}
\item Microsoft Flight Simulator -  
\item VBS (Arma) - środowisko symulacyjne do wizualizacji pola walki
\item Warthunder - darmowa gra komputerowa wprowadzająca znaczną ilość historycznych i współczesnych modeli samolotów, których parametry zostały oparte na dostępnych i odtajnych danych.
\item RealFlight - modelarski symulator lotu
\end{itemize}

Istnieją również rozwiązania typu open-source, realizujące jedynie poszczególne zadania:

\begin{itemize}
\item JSBsim - rozbudowany silnik dynamiki lotu działający w czasie rzeczywistym  
\item Ardupilot, INAV, Betaflight - kompletne systemy sterowania dla modeli zdalnie sterowanych
\end{itemize}






\section{Specyfikacja}



\subsection{Architektura}

System został podzielony na moduły które robią blabla bla...

\begin{itemize}
\item UAV\_physic\_engine 
\item UAV\_controller 
\item UAV\_visualization 
\item UAV\_drop\_physic 
\item UAV\_aggregator 
\item UAV\_server 

\end{itemize}

\newpage

\section{Dobór technologii}

\subsection{UAV\_physic\_engine}

C++ - używamy C++ ze względu na wydajność i elastyczność

Eigen - blibkioteka do algebry liniowej i macierzowej

ZeroMQ - kolejka wiadomości

RapidXML

Cxxopts

\subsection{UAV\_visualization}

Java

LWJGL

Music

\subsection{UAV\_aggregator}

Rust

Nalgebra

....


\subsection{UAV\_server}

Docker

ubuntu:latest



\color{black}
\section{Ocena ryzyka -- analiza SWOT}


\begin{tikzpicture}
\renewcommand{\arraystretch}{1.2}
\setlist{left=1em,parsep=0.5ex,after=\smallskip}
\def\myw{7.5cm}
\matrix[SWOT] 
{
\& |[fill=black!10]|\renewcommand{\arraystretch}{1.3}\begin{tabular}{Wc{\myw}Wc{\myw}}
Pozytywne & Negatywne\\ 
\end{tabular}\\
 Wewnętrzne
  \& \begin{tabular}{p{\myw}p{\myw}}
	\textbf{Silne strony:} \begin{itemize}
 	 \item przejrzysta implementacja zgodna z paradygmatami programowania obiektowego,
        	 \item modułowość projektu, pozwalająca na modyfikację poszczególnych komponentów bez konieczności przebudowy całego projektu,
        	 \item uniwersalny model dynamiki statków powietrznych, pozwalający na obliczenia w czasie rzeczywistym.
\end{itemize}
& 
 \textbf{Słabe strony:} \begin{itemize}
  	\item ograniczony czas projektu może skutkować niedopracowaniami we wdrożonych funkcjonalnościach,
        \item obliczenia symulacji i kolizji obywają się na CPU, co może ograniczać wydajność,
        \item model matematyczny jest mniej dokładny niż rozbudowana symulacja mechaniki płynów.
\end{itemize} \end{tabular}\\
 Zewnętrzne \& \begin{tabular}{p{\myw}p{\myw}}
	\textbf{Szanse:} \begin{itemize}
   	\item dzięki udostępnieniu systemu na otwartej licencji, możliwe jest wykorzystanie wypracowanych rozwiązań w przyszłych projektach,
       	\item system może stanowić narzędzie ułatwiające opracowanie i testowanie nowatorskich systemów sterowania.
\end{itemize} & \textbf{Zagrożenia:} \begin{itemize}
  \item język Rust wykorzystany w UAV\_aggregator może przestać być wspierany na przestrzeni najbliższych lat,
  \item trudności z identyfikacja wiarygodnych współczynników modelu dynamiki.
  
\end{itemize} \end{tabular}\\
};
\end{tikzpicture}






\newpage
\section{Bibliografia}

\printbibliography

%\printbibliography[type=book,heading=subbibliography,title={Literatura}]
%\printbibliography[type=article,heading=subbibliography,title={Artykuły}]
%\printbibliography[type=online,heading=subbibliography,title={Źródła internetowe}]

\end{document}\grid


\begin{comment}

Symulacje komputerowe dynamiki ruchu  W szczególności w zagadnieniu jakim jest projektowanie systemów sterowania do bezzałogowych statków powietrznych, zastosowanie takich narzędzi pozwala zminimalizować koszty wytworzenia poprawnie działającego systemu i przyśpieszyć jego rozwój i testowanie.

Celem niniejszej pracy jest opracowanie wirtualnego środowiska do symulacji dynamiki lotu bezzałogowych statków powietrznych. System implementuje podstawowy model dynamiki statków powietrznych wyposażonych w silniki rotorowe, silniki odrzutowe, powierzchnie nośne i powierzchnie sterowe. Pozwala na przeprowadzenie lotu symulowanym obiektem ktorego parametery określane są przez konfiguracje. System dzieli sie na serwer i aplikacje kliencką. Róznorodność modułów pozwala na realizacje róznych scenariusz (strzał, upuszczenie ładunku, kolizje, wpływ warunków srodowiskowych etc.)


\section{Model dynamiki statku powietrznego}

\begin{equation}
\begin{cases}
\dot{Y} =  T(Y) \cdot X + stabilizacja\\ 
A \cdot \dot{X} + \Omega (X) \cdot A \cdot X = F_g + F_a + F_d + F_o
\end{cases}
\end{equation}

gdzie:
\begin{itemize}
\item $Y$ - wektor pozycji i orientacji wyrażony w układzie globalnym
\item $X$ - wektor prędkości postępowej i kątowej w układzie związanym ze statkiem powietrznym
\item $A$ - macierz masowa
\item $\Omega$ - macierz gyroskopowa
\item $F_g$ - siła i moment pochodząca od siły grawitacji wyrażone w układzie związanym ze statkiem powietrznym
\item $F_a$ - siła i moment aerodynamiczny wyrażone w układzie związanym ze statkiem powietrznym
\item $F_d$ - siła i moment zespołów napędowych wyrażone w układzie związanym ze statkiem powietrznym
\item $F_o$ - siła i moment oddziaływań zewnętrznych napędowych wyrażone w układzie związanym ze statkiem powietrznym
\end{itemize}

\begin{equation}
Y = \begin{bmatrix}
x_{NED}\\
y_{NED}\\
z_{NED}\\
q_0\\
q_x\\
q_y\\
q_z
\end{bmatrix}
\end{equation}

\begin{equation}
X = \begin{bmatrix}
\dot{x}_b\\
\dot{y}_b\\
\dot{z}_b\\
P_b\\
Q_b\\
R_b
\end{bmatrix}
\end{equation}

\begin{equation}
F_g = R_{nb}(Y) \cdot  \begin{bmatrix}
0\\
0\\
g
\end{bmatrix}
\end{equation}

\begin{equation}
F_a = \frac{1}{2}\rho \cdot V_{tot}^2 S R_{wb} C_F \\
M_a = \frac{1}{2}\rho \cdot V_{tot}^2 S d R_{wb} C_F
\end{equation}

\end{comment}
